\documentclass{article}
\usepackage{graphicx} % Required for inserting images
\usepackage{amsmath}
\usepackage{amssymb}

\title{Lagrange Multiplier}
\author{Zhixin WANG}
\date{November 2026}

\begin{document}

\maketitle

\section{Why we need it?}
Imagine we have an optimization problem with some constraints, say $g(x) = b$, how can we find the maximum or minimum under such constraint?

Obviously, without that constraint, we can simply take derivatives (first-order + second order sometimes) to find the maximum or minimum, but the constraint part is annoying. How should we solve it?

\section{Idea}

\subsection{Lemma}
The gradient of $g$ at point $P_{0}$ is always perpendicular to the level surface $g = k$ that $P_{0} \in g = k$.

Proof:

Let's randomly pick a curve $\gamma(t)$ on the level surface $g = k$. And from the chain rule, we know:
$$\frac{\mathrm{d}}{\mathrm{d}t} g(\gamma(t)) = \nabla g(\gamma(t)) \cdot \gamma'(t) = \frac{\mathrm{d}}{\mathrm{d}t} k = 0$$
So, obviously, $\nabla g(\gamma(t)) \cdot \gamma'(t) = 0$, which means that the gradient of $g$ is perpendicular to every curve on the level surface. That is: the gradient is always perpendicular to the level surface.

\subsection{Corollary}
In the first derivative test, for $P_{0}$ to be a local min/max, the restriction of the function to all curves $\gamma(t)$ through $P_{0}$ still has a local min / max at $t = 0$. Then, we use the usual first derivative test for single variable functions and get:
$$D_{u}f (P_0) = 0 \text{ for all vector } u$$

Note that without any constraint, our curve can go anywhere and hence, the vector $u$ (in the directional derivative) can be anything. Now, subject to the constraint $g(x, y, z) = k$, we can only choose curves on the level-surface. Thus, the $u$ that appears in
$$Df(P_0) \cdot u = D_{u}f(P_0) = 0$$
can only be tangent to the level surface $g(x, y, z) = k$ at that point $P_0$. Thus, if $P_0$
is a local min/max (subject to the constraint g(x, y, z) = k), $D f (P_0)$ is perpendicular to the tangent plane to the level surface at $P_0$ and plus $Dg(P_0)$ is also perpendicular to this tangent plane (from our lemma). Thus, 
$$D f (P_0) = \lambda Dg(P_0) \text{ for some } \lambda \in \mathbb{R}$$
This is the analog of the first derivative test.

\section{Theory}
A point $P_0$ on the level surface $g(x, y, z) = k$ is a critical point of a
function $f$ restricted to the level surface if 
$$Df (P_0) = \lambda Dg(P_0) \text{for some} \lambda \in \mathbb{R}$$

Coming back to the original question: to find min/max of $f$ subject to the constraint $g(x, y, z) = k$, we look for all critical points, i.e. $P_0$ on the level surface
$g(x, y, z) = k$ such that
$$Df (P_0) = \lambda Dg(P_0) \text{for some} \lambda \in \mathbb{R}$$
Then, plug these values back to $f$ and compare the results to find min/max.


\end{document}